% Source: https://artofproblemsolving.com/community/c3h1870915_help_me_find_an_easier_way_to_solve_this_please
% Stable identity: aops-topic-1870915-q01
Count the $3$s position by position.
\begin{itemize}
\item \textbf{Units digit.} The numbers $3,13,23,\dots,573$ end in $3$: there are $\dfrac{573-3}{10}+1=58$ of them.
\item \textbf{Tens digit.} The blocks $30$--$39$, $130$--$139$, $230$--$239$, $330$--$339$, $430$--$439$, $530$--$539$ each give $10$ numbers: $60$ in total ($576$ is past $539$).
\item \textbf{Hundreds digit.} The numbers $300$--$399$: $100$ of them.
\end{itemize}
A number with several $3$s is counted once for each $3$, which is what we want. The total is $58+60+100=\boxed{218}$.
